% !TeX root = ../Rules.tex
% !TeX spellcheck = en_US
\section{Forbidden Actions and Penalties}
\label{sec:forbidden_actions}

The following actions are forbidden. In general, when a penalty applies, the robot shall be removed, not the ball.

\subsection{Penalty Procedure}
\label{sec:penalty_procedure}

When a robot commits an infraction, the head referee shall call out the \textbf{infraction} committed, the \textbf{team color} of the robot, and the \textbf{jersey number} of the robot. The penalty for the infraction will be applied immediately. The robot handlers should perform the actual movement of the robots for the penalty so that the head referee can continue focusing on the game. The operator of the GameController will send the appropriate signal to the robots indicating the infraction committed.

If a robot commits an offense in the vicinity of the ball (within the free kick avoidance radius, which is the center circle radius for the respective field size—see Table~\ref{tab:field_dim}), the game is interrupted and the offending player is removed according to the standard removal penalty. After the offending robot is removed, the opposing team restarts play with a free kick (\cf~\cref{sec:free_kick}). If the offense occurred away from the ball, the offending robot is penalized and removed, but play is not interrupted.

For penalties that are timed, the penalty time is considered to be over at the end of each half.

\subsection{Standard Removal Penalty}
\label{sec:removal_penalty}

Unless otherwise stated, all infractions result in the removal of the infringing robot from the field of play for a particular amount of time, after which it will be returned to the field of play. This process is called the \textit{standard removal penalty}.

When the head referee indicates an infraction has been committed that results in the standard removal penalty, the robot handler for that team will remove the robot immediately from the field of play. The robot should be removed in such a way as to minimize the movement of the other robots and the ball. If the ball is inadvertently moved when removing the robot, the ball should be replaced to the position it was in when the robot was removed.

The GameController will send the appropriate penalty signal to the robot indicating the infraction committed. After a penalty is signaled to the robot, it is not allowed to move in any fashion (except for standing up).

The initial duration of the standard removal penalty time is \qty{\StandardPenaltyTime}{\second}. Unless otherwise specified, the penalty time increases by \qty{\StandardPenaltyIncrease}{\second} each time a team commits any infraction. During the \texttt{set} state the penalty time counter will not decrease. The GameController will keep track of the time of the penalty.

The removed robot will be placed outside the field at a distance of approximately \qty{50}{\centi\metre} away from the nearest touchline, facing \textbf{towards} the field of play. It should be placed close to the position where the penalty mark projects onto the touchline in the robot's own half. If there is another robot already in this position, the robot should be placed at a nearby location along the touchline. When finding a nearby location, it \textbf{must} still be in the robot's own half, so that the symmetry of the field can be resolved by the robot's localization system.

After the penalty time elapses, the GameController operator will remove the penalty from the penalized robot, enabling it to autonomously re-enter the game without requiring additional intervention from the referees.

\subsection{Forbidden Actions}

The following actions are forbidden but may not result in removal penalties. Each forbidden action specifies the actions to be taken by the referees.

\subsubsection{Manual Interaction by Team Members}

Manual interaction with the robots, either directly or via some communications mechanism, is not permitted. Team members (that are not robot handlers) can only touch one of their robots when a robot handler hands it over to them after a ``Request for Pick-up'' (\cf~\cref{sec:request_for_pick_up}).

\subsubsection{Damage to the Field}
\label{sec:damage}

A robot that damages the field, or poses a threat to spectator safety, will be removed from the field for the remainder of the game.

\subsection{Illegal Positioning}
\label{sec:illegal_positioning}

A robot penalized under illegal positioning has the ``Illegal Position'' penalty applied. Illegally positioned robots are subject to the standard removal penalty (\cf~\cref{sec:removal_penalty}). The head referee will call ``Illegal Position \textless robot color\textgreater \textless robot number\textgreater''. Referees may use ``Illegal Defender'' interchangeably to help with clarity.

For simplicity, Illegal Positioning penalties during the \texttt{set} state (for kick-off or a penalty kick) do not count towards the incremental penalty count.

\subsubsection{Before and During Kick-off}
\label{sec:ip_kick_off}

If a robot violates the positioning restrictions for kick-off (\cf~\cref{sec:kick_off}) at the time the \texttt{set} state starts, it will be penalized and removed for \qty{15}{\second}. If a robot enters the areas restricted to the other team before the ball is in play, it will be penalized and removed according to the standard removal penalty. A robot making reasonable effort to return to the field from a standard penalty must not be penalized for illegal position.

\subsubsection{Own Goal Area}
\label{sec:ip_own_goal_area}

During the \texttt{set} and \texttt{playing} states, at most \textit{three} players can be within their own goal area at the same time.

A robot is within the goal area if any part of its body is touching the ground inside the goal area or touching one of its lines. The penalty is applied when any additional players enter the area in \texttt{playing}, or to the excessive players closest to the border of the goal area in \texttt{set}. Note that if a player is pushed into the goal area by an opponent, this robot will not be subject to removal, unless it fails to exit the area within \qty{5}{\second} (or \qty{5}{\second} after getting up if the pushing led to falling).

If an illegal defender kicks an own goal, the goal is scored for the opponent. If there is any doubt about whether a goal should count (e.g., the illegal defender infraction is called, but the robot scores the own goal immediately afterwards, before it is removed), then the decision shall be against the infringing robot.

\subsubsection{Defender Encroachment During Free Kick}
\label{sec:ip_free_kick}

If a robot of the defensive team enters or does not attempt to leave the circular avoidance area around the ball after a free kick (\cf~\cref{sec:free_kick}) is called, ``Illegal Position'' is called. The avoidance radius is defined as the center circle radius for the respective field size (see Table~\ref{tab:field_dim}). Note that the referee should not look for exact distances and rather penalize only those robots who clearly violate this rule. As a guideline, the robots of the defensive team should clear the ball within \qty{10}{\second}.

\subsubsection{Penalty Area During Penalty Kick}
\label{sec:ip_penalty_kick}

If a robot enters the relevant penalty area during a penalty kick, except for the goalkeeper (defending team) and one robot of the attacking team, ``Illegal Position'' is called.

\subsection{Forbidden Motion}
\label{sec:forbidden_motion}

Robots that begin moving their legs or locomote in any fashion before reaching the specified state in the following subsections will be penalized \textit{in place} on the field. ``Forbidden Motion'' penalties do not follow the standard removal procedure, and hence do not count towards the incremental penalty count. A robot will be moved back to its original position if it has moved significantly before being penalized.

\subsubsection{Motion in Set}
\label{sec:motion_in_set}

If a robot begins moving during the \texttt{set} state (i.e., before the head referee blows the whistle), it will be penalized for \qty{15}{\second} in the \texttt{playing} state. The head referee will call ``Forbidden Motion in Set \textless robot color\textgreater \textless robot number\textgreater''. Note that falsely responding to a whistle on another field will result in this penalty.

\subsection{Fallen or Inactive Robot}
\label{sec:fallen_robot}

\subsubsection{Fallen Robot}

If a robot falls during the game, it should start executing a get-up action within \qty{\FallenRobotGetUpTime}{\second}. If it does not commence a get-up action within this time, it will be penalized and removed for \qty{\FallenRobotPenaltyTime}{\second}.

A fallen robot which is unable to autonomously stand up within \FallenRobotGetUpAttempts{} attempts, or \FallenRobotGetUpAttemptsDisturbed{} attempts if disturbed externally during the first attempts, will be penalized and removed for \qty{\FallenRobotPenaltyTime}{\second}. The head referee will call ``Fallen Robot \textless robot color\textgreater \textless robot number\textgreater''.

The goalkeeper, inside its own penalty area, is the only robot permitted to ``dive'' (deliberately fall in a way that might cause its torso, arms, or hands to intercept the ball). In all other cases, robots should remain upright—that is, supported by their feet.

Fallen Robot penalties do not follow the standard removal procedure and do not count towards the incremental penalty count.

\subsubsection{Inactive Robot}

A robot that has ceased activity for \qty{\InactiveRobotTime}{\second} or has turned off will be removed and penalized for \qty{\FallenRobotPenaltyTime}{\second}. The head referee will call ``Inactive Robot \textless robot color\textgreater \textless robot number\textgreater''.

A robot is considered active if it performs at least one of the following:
\begin{enumerate}
  \item The robot walks in any direction, or turns.
  \item The robot searches for the ball, or is tracking the ball.
\end{enumerate}

Inactive Robot penalties do not follow the standard removal procedure and do not count towards the incremental penalty count.

\paragraph{Note:} The intention of this rule is not to penalize robots simply for being stationary—provided they are not ``asleep'' and have not ``crashed''.

\subsection{Local Game Stuck Penalty}
\label{sec:pen_local_game_stuck}

When Local Game Stuck is called (\cf~\cref{sec:local_game_stuck}), the nearest robot to the ball will be penalized and removed for \qty{\FallenRobotPenaltyTime}{\second}. Local Game Stuck penalties do not follow the standard removal procedure and do not count towards the incremental penalty count.

\subsection{Ball Holding}
\label{sec:ball_holding}

The goalkeeper is allowed to hold the ball for up to \qty{\BallHoldingTimeGoalkeeper}{\second} as long as it has at least one foot inside its own penalty area. In all other cases (except those noted in~\cref{sec:ball_holding_exceptions}), robots are allowed to hold the ball for up to \qty{\BallHoldingTime}{\second}. Holding the ball for longer than this is not allowed. The head referee will call ``Ball Holding \textless robot color\textgreater \textless robot number\textgreater'', and the robot is removed under the standard removal penalty.

The ball should be removed from the possession of the robot and placed where the penalty occurred. If the robot that held the ball has moved the ball before the robot can be removed, the ball shall be replaced where the penalty occurred. This applies to accidental goals.

\textbf{Example:} A robot holds the ball, and before the referees can remove the robot, it shoots the ball into the goal. The goal will not be counted, and the ball will be replaced where the penalty occurred.

A robot must leave enough open space around the ball. The occupation of the ball is judged using the convex hull of the projection of the robot's body onto the ground. ``Enough open space'' means that at least half of the ball is not covered by the convex hull. It is not important whether the robot actually touches the ball.

Intentional continual holding is prohibited even if each individual holding time does not exceed the time limit. In general, robots should release the ball for approximately as long as they were holding it to reset the clock. Without a sufficient release, the continual holding is regarded as a continuous hold from the very beginning and the holding rule is strictly applied.

\subsubsection{Exceptions to Ball Holding}
\label{sec:ball_holding_exceptions}

The following situations are exceptions where ball holding does not apply:
\begin{enumerate}
  \item Ball holding may not be called when the ball becomes stuck between a robot's legs. In such a situation, the head referee should call ``Clear Ball'' and a robot handler or assistant referee should remove the ball and place it approximately where it was before it became stuck.
  \item Ball holding may not be called when a robot falls on a ball. The robot will either get up and hence free the ball, or the robot should be removed under the Fallen Robot rule.
\end{enumerate}

\subsection{Player Stance}
\label{sec:player_stance}

Robots must not maintain a stance wider than a reasonable proportion of their body width for longer than \qty{10}{\second}. As a guideline, the stance should not exceed approximately 1.5 times the robot's shoulder width. Staying in an excessively wide stance for longer than \qty{10}{\second} will result in the standard removal penalty. \unsure{Review call phrase here, should this be an existing call in GC for simplicity?}

This rule applies to upright robots only, i.e., robots that are supported by their feet. Other cases are handled by the Fallen Robot rule (\cf~\cref{sec:fallen_robot}).

\subsection{Player Pushing}
\label{sec:player_pushing}

\emph{Pushing} is a direct or indirect forceful contact with any opponent robot, sufficient to destabilize it, and is not allowed. The head referee will call ``Pushing \textless robot color\textgreater \textless robot number\textgreater''.

If the ball moves significantly as a result of pushing, then it should be replaced to where it was at the time of the infraction.

A \textbf{Pushing Free Kick} (\cf~\cref{sec:free_kick}) is awarded against the robot penalized for pushing:
\begin{itemize}
  \item if the robot is in an approximately \qty{0.5}{\metre} radius around the ball, and
  \item if the pushed robot was \textbf{not} inside the penalty area of the pushing robot.
\end{itemize}

A \textbf{Penalty Kick} (\cf~\cref{sec:penalty_kick}) is awarded against the robot penalized for pushing:
\begin{itemize}
  \item if the robot is in an approximately \qty{0.5}{\metre} radius around the ball, and
  \item if the pushed robot was inside the penalty area of the pushing robot.
\end{itemize}

The following exceptions define situations where pushing does not apply:
\begin{itemize}
  \item Pushing may occur \textbf{only} between players of different teams.
  \item A stationary robot cannot be penalized for pushing, including a robot that is kicking, provided that the ball was close enough where a kick could have succeeded at the start of the kick motion.
  \item A robot currently getting up cannot be penalized for pushing.
  \item The goalkeeper cannot be penalized for pushing while looking at or chasing the ball in its own penalty area.
  \item Front-to-front contact between robots with the ball between them does not constitute pushing, unless one robot is moving with significantly more force or speed than the other can reasonably withstand given the size and weight differences between the robots.
  \item Any robot proceeding to the ball whose side (arm, shoulder, etc.) makes contact with another robot cannot be called for pushing, even if the second robot is not proceeding to the ball.
  \item A robot pushed by another robot cannot simultaneously be called for pushing itself. Only the robot who initially pushed the other robot will be called for pushing.
\end{itemize}

\subsection{Playing with Arms/Hands}
\label{sec:playing_with_hands}

Playing with arms/hands occurs when a field player or a goalkeeper outside its own penalty area moves its arms/hands to touch the ball (except during a fall or get-up). The goalkeeper is allowed to touch the ball with its arms/hands while it is within its own penalty area.

A robot playing with arms/hands will be subject to the standard removal penalty and the ball will be replaced at the point where it contacted the arms/hands of the offending robot. If an own goal is scored as a result of playing with arms/hands, the goal should count and the player should not be penalized.

Accidental playing with arms/hands when a robot falls or executes a get-up routine will not be penalized. If the ball goes out of play in this case, normal kick-in rules will apply (\cf~\cref{sec:kick_in}). However, goals (except for own goals) resulting from a ball contact with the arms/hands during a fall or get-up do not count and result in a Goal Kick (\cf~\cref{sec:goal_kick}) as if the ball went over the goal line next to the goal.

\subsection{Leaving the Field}
\label{sec:leaving_the_field}

A robot that intentionally leaves the carpeted playing area will be subject to the standard removal penalty (\cf~\cref{sec:removal_penalty}). The head referee will call ``Leaving the Field \textless robot color\textgreater \textless robot number\textgreater''.

Additionally, a robot will also be subject to the standard removal penalty when:
\begin{itemize}
  \item The robot walks into the goalposts or goal net for more than \qty{5}{\second}, including robots that are stuck on the goalposts and unable to free themselves.
  \item The robot's fingers or other appendages become entangled in the net (without any time constraint).
\end{itemize}

\subsection{Jamming}
\label{sec:jamming}

During a match, robots shall not jam the communication and sensor systems of the opponents:

\begin{description}
  \item[Wireless communication.] Each team is allowed to send a maximum of \UDPPacketLimit{} UDP packets of up to \qty{\UDPPacketSize}{\byte} per game, in addition to \UDPDebugPacketRate{} packet per second per robot to a single team device on the wired network for debugging purposes (no size constraint for debug packets, but must be within the maximum single UDP packet size). Any team that exceeds the team communication limits, as tracked by the GameController, will have all goals automatically nullified for the game. The GameController will communicate at \GameControllerPacketRate{} packet(s) per second. If a team violates these limits, penalties will apply. If a team violates this rule in multiple games, disqualification from the tournament (including all technical challenges and side competitions) will be a potential consequence. Except for the wireless equipment and access points provided by the organizers of the competition, nobody close to the field is allowed to use \qty{2.4}{\giga\hertz} or \qty{5}{\giga\hertz} radio equipment (including cellular phones and/or Bluetooth devices).
  
  \item[Whistle interference.] Both teams and the audience shall avoid intentionally confusing the robots by producing sounds similar to the game whistle.
  
  \item[Acoustic communication.] If acoustic communication is used by both teams, they shall negotiate before the match how they can reduce interference. If only one team uses acoustic communication, the robots of the other team shall avoid producing any sound. In addition, both teams and the audience shall avoid intentionally confusing the robots by producing similar sounds to those used for communication.
  
  \item[Visual perception.] The use of flashlights or other bright light sources is not allowed during games. However, flash photography from the audience is allowable as long as the head referee believes the purpose is not to jam any of the robots.
\end{description}

\subsection{Aborting a Penalty}
\label{sec:aborting_penalty}

If a penalty was incorrectly applied to a robot and acknowledged by the referee to be incorrect, the robot should first be removed and placed outside the field, following the same process as for the Standard Removal Penalty (\cf~\cref{sec:removal_penalty}). Once placed, the robot should then be allowed to continue play immediately after being placed and the GameController removes its penalty.

The aborted penalty will not be added to the total team penalty count and will not contribute to the incremental penalty increase.

\subsection{Request for Pick-up}
\label{sec:request_for_pick_up}

A robot handler may request to pick up a robot if and only if a robot is in a dangerous situation that is likely to lead to physical injuries. The head referee must approve this request prior to the robot handler touching the robot. If a team member touches a robot without the permission of the referee, the respective robot and the person touching it each receive an official warning.

The referee may deny a request for pick-up. A pick-up is not granted for reasons such as being incorrectly configured or being unable to see the ball.

\subsection{Untangling}
\label{sec:untangling}

If entangled robots fail to untangle themselves, the referee may ask designated robot handlers of both teams to untangle the robots. Untangling must not make significant changes to robot positions or heading directions. Untangled robots must be placed on the ground not closer than \qty{50}{\centi\metre} to the ball and in a way that does not give either team an advantage. Being entangled does not incur a penalty.

\subsection{Serious Foul Play}
\label{sec:serious_foul_play}

A player who commits serious foul play must be removed from the game for the remainder of the match. Serious foul play includes:

\begin{itemize}
  \item Endangering any human
  \item Repeated deliberate violations of the rules
  \item Actions that could cause significant damage to other robots
\end{itemize}

A player removed for serious foul play may not be substituted. As a result, the team must play the remainder of the match with one fewer player.

\subsection{Substituting a Penalized Player}
\label{sec:substituting_a_penalized_player}

Teams may substitute a player who has received official warnings as usual. However, all warnings issued to the substituted player are transferred to the player entering the field.

\subsection{Penalties against Teams or Humans}
\label{sec:penalties_against_humans}

The referee can issue yellow and red cards against humans or teams at their discretion. Cards can be issued for offenses such as:

\begin{itemize}
  \item Offensive or abusive behavior
  \item Repeatedly ignoring referee instructions
  \item Repeatedly entering the field without permission
  \item Breaking the game rules to gain an advantage
\end{itemize}

Cards that are issued against humans or teams must be reported to the Technical Committee immediately after the game. Cards against humans or teams persist between games. A human who receives multiple cards may be prohibited from entering the field or team area during games for the remainder of the competition. A team that commits serious or repeated violations may be disqualified from the competition. 